\documentclass[11pt,a4paper]{article}

% === 中文支持 ===
\usepackage{ctex}
\usepackage[T1]{fontenc}

% === 页面布局 ===
\usepackage[top=2.5cm, bottom=2.5cm, left=2.5cm, right=2.5cm]{geometry}

% === 表格 ===
\usepackage{booktabs}
\usepackage{array}
\usepackage{enumitem}

% === 数学符号 ===
\usepackage{amssymb}

% === 代码 ===
\usepackage{listings}
\lstset{basicstyle=\ttfamily\footnotesize, breaklines=true, frame=single, showstringspaces=false}

% === 颜色 ===
\usepackage{xcolor}
\definecolor{primary}{HTML}{1a1a2e}
\definecolor{accent}{HTML}{3b82f6}
\definecolor{success}{HTML}{22c55e}
\definecolor{danger}{HTML}{ef4444}
\definecolor{warning}{HTML}{f59e0b}
\definecolor{graytext}{HTML}{6b7280}
\definecolor{progress}{HTML}{8b5cf6}

% === 超链接 ===
\usepackage[hidelinks]{hyperref}
\hypersetup{colorlinks=true, linkcolor=primary, urlcolor=accent}

% === 标题样式 ===
\usepackage{titlesec}
\titleformat{\section}{\Large\bfseries\color{primary}}{}{0em}{}[\titlerule]
\titleformat{\subsection}{\large\bfseries\color{primary}}{}{0em}{}

% === 页眉页脚 ===
\usepackage{fancyhdr}
\pagestyle{fancy}
\fancyhf{}
\fancyhead[L]{\small\color{graytext} 企业级 AI Agent 应用商店 -- 需求文档}
\fancyhead[R]{\small\color{graytext} DEMAND-002}
\fancyfoot[C]{\small\color{graytext} \thepage}
\renewcommand{\headrulewidth}{0.4pt}

% === 文档信息 ===
\title{\textbf{DEMAND-002：开发环境搭建与依赖安装}\\[0.5em]\small V1 阶段首个需求}
\author{万能产品小助手（PM AI）\\魏子政 审阅}
\date{2026-05-06 / 版本 v0.1}

\begin{document}

\maketitle

\begin{center}
\begin{tabular}{>{\bfseries}p{3cm} p{10cm}}
\toprule
提出日期 & 2026-05-06 \\
提出者 & 魏子政 \\
优先级 & \textcolor{danger}{P0} \\
状态 & \textcolor{progress}{进行中} \\
阶段 & V1（正式需求阶段） \\
总需求编号 & D-008 \\
影响范围 & 项目根目录结构、Docker Compose 配置、前后端项目初始化 \\
\bottomrule
\end{tabular}
\end{center}

\vspace{1em}

\section{需求描述}

搭建完整的开发环境，确保前后端项目可运行、数据库/搜索/存储服务可连接。通过 \texttt{docker compose up} 一键启动所有依赖服务，Cursor 可以直接在项目目录下开始功能开发。

这是 V1 阶段的第一步，后续所有功能开发（CRUD、搜索、审批、安全扫描）都依赖本需求的完成。

\section{背景与动机}

V0 先备阶段已完成技术选型、智能体分工、文档体系等基础工作。但项目尚未有任何可运行的代码。需要先打通「环境搭建 $\to$ 服务启动 $\to$ 健康检查」这条链路，验证技术选型的可行性。

\section{具体任务}

\subsection{1. 后端项目初始化}

\begin{itemize}[nosep]
  \item 创建 Python 项目结构（FastAPI）
  \item 安装核心依赖：fastapi, uvicorn, sqlalchemy[asyncio], asyncpg, pydantic, python-jose, passlib, alembic
  \item 配置 SQLAlchemy async engine + session
  \item 配置 Alembic 数据库迁移
  \item 编写 \texttt{/health} 健康检查接口
  \item 配置 CORS（允许前端跨域）
\end{itemize}

\subsection{2. 前端项目初始化}

\begin{itemize}[nosep]
  \item 使用 Vite 创建 Vue 3 + TypeScript 项目
  \item 安装 Element Plus + Pinia + Vue Router
  \item 配置 Vite 开发服务器代理（\texttt{proxy} 指向后端 API）
  \item 创建基础页面框架（Layout + 首页）
\end{itemize}

\subsection{3. Docker Compose 编排}

\begin{itemize}[nosep]
  \item PostgreSQL 16（端口 5432）
  \item MinIO（API 端口 9000，控制台端口 9001）
  \item OpenSearch 2.x（端口 9200，单节点模式）
  \item 后端 FastAPI 服务（端口 8000）
  \item 网络和数据卷配置
  \item \texttt{.env} 环境变量文件（数据库连接、密钥等）
\end{itemize}

\subsection{4. 环境验证}

\begin{itemize}[nosep]
  \item \texttt{docker compose up -d} 启动所有服务无报错
  \item \texttt{curl http://localhost:8000/health} 返回 200
  \item 前端 \texttt{http://localhost:5173} 可访问
  \item \texttt{alembic upgrade head} 执行成功
  \item MinIO 控制台 \texttt{http://localhost:9001} 可访问
\end{itemize}

\section{预期目录结构}

\begin{lstlisting}[basicstyle=\ttfamily\footnotesize, frame=single]
enterprise-app-store/
  docker-compose.yml
  .env.example
  backend/
    app/
      main.py              # FastAPI 入口
      config.py            # 配置
      database.py          # SQLAlchemy async
      models/              # ORM 模型
      api/                 # 路由
    alembic/               # 迁移文件
    alembic.ini
    requirements.txt
  frontend/
    src/
      App.vue
      main.ts
      router/
      stores/
      views/
    package.json
    vite.config.ts
\end{lstlisting}

\section{验收标准}

\begin{itemize}
  \item[\checkmark] \texttt{pip install -r backend/requirements.txt} 全部成功
  \item[\checkmark] \texttt{cd frontend \&\& npm install} 全部成功
  \item[\checkmark] \texttt{curl http://localhost:8000/health} 返回 \texttt{\{\"status\":\"ok\",\"version\":\"0.1.0\"\}} (200 OK)
  \item[\checkmark] \texttt{npx vite build} 构建成功（3.59s）
  \item[\checkmark] docker-compose.yml 包含 postgres/minio/opensearch/backend 四个服务
  \item[\checkmark] .env.example 环境变量完整
  \item[\checkmark] \texttt{curl http://localhost:8000/health} returns status=ok (200 OK)
  \item[$\square$] \texttt{docker compose up -d} 启动所有服务（待 Docker 环境验证）
  \item[$\square$] \texttt{alembic upgrade head} 数据库迁移（待 Docker 环境验证）
\end{itemize}

\vspace{1em}
\textit{本地验证已通过（2026-05-06）。Docker 环境验证待后续。}

\section{分配方式}

由 PM AI（OpenClaw）调用 Cursor Agent 执行：

\begin{lstlisting}[basicstyle=\ttfamily\footnotesize, frame=single]
cursor agent --print --trust \
  --workspace projects/enterprise-app-store \
  "搭建开发环境：1.初始化后端FastAPI项目 2.初始化前端Vue3项目
   3.编写docker-compose.yml 4.配置.env 5.验证所有服务启动正常"
\end{lstlisting}

\section*{更新记录}

\begin{tabular}{lll}
\toprule
版本 & 日期 & 变更内容 \\
\midrule
v0.1 & 2026-05-06 & 初始版本 \\
\bottomrule
\end{tabular}

\end{document}
